\section{ویژگی‌های جت‌ها و ذره‌ها}
\label{sec:vars}
برای هر جت یک مجموعه ویژگی زیرساختاری و سینماتیکی، و برای هر ذره داخل جت (تا ۱۰۰ ذره) ۱۷ ویژگی ذخیره می‌شود. در این بخش جزئیات این ویژگی‌ها را توضیح می‌دهیم.

\subsection{ویژگی‌های جت‌ها}
اولین ویژگی‌ها سینماتیکی‌اند. تکانه عرضی ($\pt$) مؤلفه تکانه عمود بر محور باریکه است، $\pt=\sqrt{p_x^2+p_y^2}$. شبه‌سرعت ($\eta$) از رابطه $\eta=-\ln\tan(\theta/2)$ به دست می‌آید که $\theta$ زاویه با محور باریکه است؛ مقدار صفر یعنی راستای عمود بر باریکه و مقدار بزرگ یعنی نزدیک به باریکه. زاویه سمتی ($\varphi$) زاویه در صفحه عمود بر باریکه است و فاصله زاویه‌ای ($\Delta R$) فاصله دو جسم در صفحه $(\eta,\varphi)$ را می‌دهد، $\Delta R=\sqrt{\Delta\eta^2+\Delta\varphi^2}$. جرم جت ($m$) جرم ناوردای مجموع چهارتکانه‌های همه ذرات تشکیل‌دهندهست، $m^2=\left(\sum_i p_i\right)^2$، و برای \Zbb قله‌ای نزدیک جرم \Zbos{} دارد.

پهنای جت ($g$) مجموع فاصله ذرات تشکیل‌دهنده از محور جت با وزن $\pt$ و نرمال‌شده با $\pt$ جت است، $g=\frac{1}{\pt^{\mathrm{jet}}}\sum_i \pt^{\,i}\,\Delta R_{i}$، و نشان می‌دهد انرژی جت چقدر از محور دور پخش شده است. ویژگی‌های مربوط به ساختار درونی جت ($\tau_1$، $\tau_2$، $\tau_3$) از رابطه $\tau_N=\frac{1}{d_0}\sum_i \pt^{\,i}\min\left(\Delta R_{i1},\dots,\Delta R_{iN}\right)$ با $d_0=R\sum_i\pt^{\,i}$ به دست می‌آیند که $\Delta R_{ik}$ فاصله ذره $i$ تا $k$-امین هسته جت است؛ مقدار کوچک $\tau_N$ یعنی جت با $N$ هسته سازگار است. نسبت‌های $\tau_{21}=\tau_2/\tau_1$ و $\tau_{32}=\tau_3/\tau_2$ وابستگی کمی به مقیاس کلی $\tau$ دارند: $\tau_{21}$ کوچک یعنی ساختار دو‌هسته‌ای (مانند \Zbb) و $\tau_{32}$ دو‌هسته‌ای را از سه‌هسته‌ای جدا می‌کند.

توابع همبستگی انرژی ($C_2$، $D_2$، $N_2$) بر پایه $e_2=\sum_{i<j} z_i z_j\,\Delta R_{ij}^{\beta}$ و $e_3=\sum_{i<j<k} z_i z_j z_k\,(\Delta R_{ij}\Delta R_{ik}\Delta R_{jk})^{\beta}$ ساخته می‌شوند که $z_i=\pt^{\,i}/\pt^{\mathrm{jet}}$ است؛ برای نمونه $C_2=e_3/e_2^2$ و $D_2=e_3/e_2^3$، و مانند $\tau_{21}$ ساختار دو‌هسته‌ای را می‌سنجند. در پایان، چند ویژگی از حذف ذره‌های نرم (\lr{soft-drop}) می‌آید که در آن ذره‌های نرمِ دور از هسته جت، یعنی آنهایی که سهم تکانه‌شان از جت کمتر از آستانه است، حذف می‌شوند. جرم پس از \lr{soft-drop} ($m_\mathrm{sd}$) جرم جت پس از این حذف است. سهم تکانه زیرجت نرم‌تر ($z_g$) برابر $z_g=\min(\pt^{(1)},\pt^{(2)})/(\pt^{(1)}+\pt^{(2)})$ برای دو زیرجت باقی‌مانده است و فاصله دو زیرجت ($R_g$) فاصله زاویه‌ای $\Delta R$ بین همان دو زیرجت است.

\subsection{ویژگی‌های ذره‌ها}
برای هر ذره داخل جت ۱۷ ویژگی ذخیره می‌شود که در آنها $\eta_\mathrm{jet}$ و $\varphi_\mathrm{jet}$ محور جت‌اند. مؤلفه‌های تکانه ($p_x$، $p_y$، $p_z$) و انرژی ($E$) ذره مستقیماً ذخیره می‌شوند و انرژی نسبی ($E_\mathrm{rel}$) برابر $E/E_\mathrm{jet}$ است. تکانه عرضی ($\pt$) و تکانه عرضی نسبی ($\pt^\mathrm{rel}=\pt/\pt^\mathrm{jet}$) نیز ذخیره می‌شوند. شبه‌سرعت ($\eta$) همراه با شبه‌سرعت نسبی ($\eta_\mathrm{rel}=\eta-\eta_\mathrm{jet}$) و زاویه سمتی ($\varphi$) همراه با زاویه سمتی نسبی ($\varphi_\mathrm{rel}=\varphi-\varphi_\mathrm{jet}$) ذخیره می‌شوند، که نسبی یعنی نسبت به محور جت. شبه‌سرعت چرخیده ($\eta_\mathrm{rot}$) و زاویه سمتی چرخیده ($\varphi_\mathrm{rot}$) به محور جت وابسته نیستند و در قابی تعریف می‌شوند که در ادامه شرح داده شده است. فاصله از محور جت ($\Delta R=\sqrt{\eta_\mathrm{rel}^2+\varphi_\mathrm{rel}^2}$)، کسینوس زاویه قطبی ذره ($\cos\theta=p_z/|\vec p|$، زاویه نسبت به محور باریکه) و کسینوس زاویه بین ذره و تکانه جت ($\cos\theta_\mathrm{rel}$) سه ویژگی دیگرند. آخرین ویژگی نوع ذره (\lr{pdgid}) است؛ برای ردهای باردار همان کد ذره در \lr{Delphes} است و برای اثرات آشکارشده در کالریمترها ۲۲ فوتون و ۱۳۰ هادرون خنثی را نشان می‌دهد. کدهای دیده‌شده ۱۱ (الکترون)، ۱۳ (میون)، ۲۲، ۱۳۰، ۲۱۱ (پیون باردار)، ۳۲۱ (کائون) و ۲۲۱۲ (پروتون) هستند.

برای ساختن قاب چرخیده، ابتدا مبدأ صفحه $(\eta,\varphi)$ را روی سخت‌ترین ذره جت می‌گذاریم، یعنی $\eta$ و $\varphi$ این ذره را از همه ذره‌ها کم می‌کنیم ($\varphi$ با در نظر گرفتن دوره‌ای بودن). سپس کل صفحه را آن‌قدر می‌چرخانیم که دومین ذره سخت روی محور منفی $\varphi$ بیفتد؛ زاویه چرخش از جهت همین ذره نسبت به محور $\eta$ به دست می‌آید. $\eta_\mathrm{rot}$ و $\varphi_\mathrm{rot}$ مؤلفه‌های هر ذره پس از این انتقال و چرخش‌اند. به همین دلیل $\eta_\mathrm{rot}$ و $\varphi_\mathrm{rot}$ ذره اول همیشه صفرند و ذره دوم فقط یک مؤلفه ناصفر دارد ($\varphi_\mathrm{rot}$ منفی)؛ از ذره سوم به بعد توزیع‌ها شکل ساختار دو‌هسته‌ای را نشان می‌دهند.

علاوه بر ویژگی‌ها، برای هر جت یک تصویر ۳۲$\times$۳۲ با وزن $\pt$ در صفحه $(\eta_\mathrm{rel},\varphi_\mathrm{rel})$ و بازه $\pm0.8$ ذخیره می‌شود. همه این‌ها در چهار فایل \lr{HDF5} با نام‌های \lr{events}، \lr{jet\_features}، \lr{constituents} و \lr{images} نگه داشته می‌شوند و پیوند رویداد$\to$جت$\to$ذره در آنها حفظ شده است.
